% Zone „Kennst du schon“ – Daten, Kl. 7
\einheitenkopf[zone][Kennst du schon]{Das kennst du schon}
\setcounter{aufgabe}{0}

\begin{aufgabe}{Ich kann den Prozentsatz berechnen.}
Wie viel Prozent sind das?
\begin{teilezwei}
\tz 3 von 10: \leerfeld[\%] & \tz 1 von 4: \leerfeld[\%] \\
\tz 9 von 20: \leerfeld[\%] & \tz 7 von 8: \leerfeld[\%] \\
\anweisung{Rechne mit dem Taschenrechner. Runde auf eine Stelle nach dem Komma.}
\tz 5 von 12: \leerfeld[\%] & \\
\end{teilezwei}
\end{aufgabe}

\begin{aufgabe}{Ich kann Prozentangaben als Dezimalzahl schreiben.}
Schreibe als Dezimalzahl.
\begin{teilezwei}
\tz $50\,\%$ = \leerfeld & \tz $25\,\%$ = \leerfeld \\
\tz $12{,}5\,\%$ = \leerfeld & \tz $5\,\%$ = \leerfeld \\
\end{teilezwei}
\end{aufgabe}

\begin{aufgabe}{Ich kann die relative Häufigkeit angeben.}
In der 7b wurde gefragt: Wie kommst du zur Schule?

\sachtabelle{lc}{Schulweg & Strichliste}{Rad & \strichliste{8}\\ Bus & \strichliste{6}\\ zu Fuß & \strichliste{4}\\ Auto & \strichliste{2}}

Zähle aus.
\begin{teile}
\teil Wie viele Kinder kommen mit dem Rad? \hfill \leerfeld
\teil Wie viele Kinder wurden insgesamt gefragt? \hfill \leerfeld
\anweisung{Gib die relative Häufigkeit als Bruch und in Prozent an.}
\teil Rad in der 7b \hfill \leerfeld \quad \leerfeld[\%]
\teil Von 400 befragten Jugendlichen fahren 14 mit dem Roller. \hfill \leerfeld \quad \leerfeld[\%]
\teil 12 Kinder sagen „ja“, 36 Kinder sagen „nein“. Anteil „ja“? \hfill \leerfeld \quad \leerfeld[\%]
\end{teile}
\end{aufgabe}

\begin{aufgabe}{Ich kann Werte aus einem Säulendiagramm ablesen.}
Eine Bäckerei zählt die verkauften Brötchen.

\saeulenab[ymax=5,ystep=1,yfein=0.25,ylabel=Brötchen in Tausend]{Mo/4,Di/3,Mi/2.75,Do/3.5,Fr/4.5}

Lies ab.
\begin{teilezwei}
\tz Montag: \leerfeld[Tausend] & \tz Dienstag: \leerfeld[Tausend] \\
\tz Mittwoch: \leerfeld[Tausend] & \\
\end{teilezwei}
\begin{teile}
\teil Wie viele Brötchen verkauft die Bäckerei am Donnerstag? \hfill \leerfeld[Brötchen]
\end{teile}
\end{aufgabe}

\begin{aufgabe}{Ich kann einen Winkel als Anteil vom Vollkreis angeben.}
Welcher Anteil vom Vollkreis (360°) ist der Winkel? Gib ihn als Bruch an.
\begin{teilezwei}
\tz 180°: \leerfeld & \tz 90°: \leerfeld \\
\tz 120°: \leerfeld & \tz 45°: \leerfeld \\
\end{teilezwei}

\anweisung{Wie groß ist der Winkel?}
\begin{teile}
\teil ein Fünftel des Vollkreises: \leerfeld[°]
\end{teile}
\end{aufgabe}

\begin{aufgabe}{Ich kann einen Winkel mit dem Geodreieck zeichnen.}
Trage den Winkel am Strahl an.

\noindent\begin{minipage}[t]{0.48\linewidth}
\begin{teile}\teil 90°\end{teile}
\winkelstrahl[6]
\end{minipage}\hfill
\begin{minipage}[t]{0.48\linewidth}
\begin{teile}\teil 60°\end{teile}
\winkelstrahl[6]
\end{minipage}

\noindent\begin{minipage}[t]{0.48\linewidth}
\begin{teile}\teil 35°\end{teile}
\winkelstrahl[6]
\end{minipage}\hfill
\begin{minipage}[t]{0.48\linewidth}
\begin{teile}\teil 140°\end{teile}
\winkelstrahl[6]
\end{minipage}
\end{aufgabe}

\begin{aufgabe}{Ich kann Längen in Zentimetern und Millimetern abtragen.}
Rechne um.
\begin{teilezwei}
\tz 3 cm = \leerfeld[mm] & \tz 50 mm = \leerfeld[cm] \\
\tz 4,5 cm = \leerfeld[mm] & \tz 7 mm = \leerfeld[cm] \\
\anweisung{Zeichne eine Strecke mit dieser Länge.}
\end{teilezwei}
\begin{teile}
\teil 6 cm \par\vspace{10mm}
\teil 3,5 cm \par\vspace{10mm}
\end{teile}
\end{aufgabe}

\begin{aufgabe}{Ich kann Kenngrößen einer Liste bestimmen.}
Punkte bei einem Quiz: 12, 7, 15, 9, 11. Bestimme.
\begin{teilezwei}
\tz kleinster Wert: \leerfeld & \tz größter Wert: \leerfeld \\
\tz Spannweite: \leerfeld & \tz Median: \leerfeld \\
\anweisung{Neue Liste: 8, 3, 10, 5, 6, 12.}
\tz Median: \leerfeld & \\
\end{teilezwei}
\end{aufgabe}

\begin{aufgabe}{Ich kann einen Bruchteil und ein Vielfaches einer Zahl berechnen.}
Berechne.
\begin{teilezwei}
\tz die Hälfte von 80: \leerfeld & \tz das Doppelte von 35: \leerfeld \\
\tz zwei Drittel von 150: \leerfeld & \tz ein Drittel mehr als 120: \leerfeld \\
\tz ein Viertel weniger als 200: \leerfeld & \\
\end{teilezwei}
\end{aufgabe}

\begin{aufgabe}{Ich finde den Fehler beim Anteil.}
Von 25 Kindern haben 10 ein Haustier, 15 haben keins. Tom berechnet den Anteil der Kinder mit Haustier:
\rechnung{\text{Anteil} &= 10 : 15 \\ &\approx 0{,}67 = 67\,\%}
Finde den Fehler, benenne ihn und korrigiere die Rechnung.
\par\vspace{22mm}
\end{aufgabe}

\begin{aufgabe}{Ich kann einen Anteil an allen berechnen.}
Wie viel Prozent sind das?
\begin{teile}
\teil Von 40 Kindern spielen 14 ein Instrument, die anderen nicht. Anteil mit Instrument?\\ \leerfeld[\%]
\teil 18 Kinder sagen „ja“, 42 Kinder sagen „nein“. Anteil „ja“? \hfill \leerfeld[\%]
\end{teile}
\end{aufgabe}
